
\begin{claim}{AC.33}{Proved; a Frobenius obstruction with a split coefficient algebra}
Let \(C/k\) have genus two in characteristic two, with
\(P\in C(k)\), and let \(\ell\ge5\) be an odd prime.
Write \(F\) for arithmetic Frobenius on \(J_C[\ell](\overline k)\).
Suppose that, for some \(d>0\),
\[
 F^d=sI,\qquad s\in\mathbf F_\ell^\times
       \text{ of order }\ell-1,
\]
and that every nonzero quotient functional has exact Frobenius
period \(d(\ell-1)\).
Then no \(P\)-split constant cyclic cover of degree \(\ell\)
satisfies the rational criterion over any finite extension of \(k\).
No assumption that two be primitive modulo \(\ell\) is needed.

In particular this conclusion holds if the characteristic
polynomial of \(F\) is irreducible over \(\mathbf F_\ell\)
and \(q=\#k\) generates \(\mathbf F_\ell^\times\).
\end{claim}
\begin{proof}
The locally exact theta divisor is defined over \(k\).
A nonzero point \(a\) of order \(\ell\) on it would, by the
scalar Frobenius identity, put every nonzero multiple of \(a\)
on theta. In genus two theta is a translate of the curve,
so the difference bound proved in AC.32 would give
\[
                    (\ell-1)(\ell-2)\le2(\ell-1),
\]
which is impossible for \(\ell\ge5\).
Thus theta has no nonzero \(\ell\)-torsion, and every geometric
cyclic \(\ell\)-cover has ordinary new factor, of two-rank
\(\ell-1\).

For any such geometric cover, its kernel is stable under
\(F^d\), and the rational point \(P\) gives its unique pointed
descent to the degree-\(d\) extension \(k_d\).
Its geometric new two-torsion \(W\) is free of rank one over
the finite reduced algebra
\[
              R=\mathbf F_2[X]/\Phi_\ell(X).
\]
Indeed each nontrivial irreducible deck-character component
occurs with multiplicity one when the new factor is ordinary.
The algebra \(R\) may be a product of several fields.
Frobenius over \(k_d\) acts on it by \(X\mapsto X^s\)
(or its inverse with dual character conventions).
This automorphism \(\sigma\) has order \(\ell-1\) and
\[
                              R^\sigma=\mathbf F_2.
\]
To check this last equality, extend scalars to
\(\overline{\mathbf F}_2\), where \(R\) is the algebra of
functions on the nontrivial characters; the primitive scalar
permutes all of them transitively, so only constants remain.
Descent of a constant function to \(R\) requires its value to
lie in \(\mathbf F_2\).

Choose an \(R\)-basis of \(W\). Frobenius is then
\(w\mapsto a\,\sigma(w)\) with \(a\in R^\times\).
Its \((\ell-1)\)-st power is multiplication by
\[
                 a\,\sigma(a)\cdots\sigma^{\ell-2}(a)
                         \in(R^\sigma)^\times=\{1\}.
\]
Therefore all new two-torsion is rational over
\(k_{d(\ell-1)}\). By the exact period hypothesis this is the
first field with constant deck group. The unique pointed
descent identifies every eligible cover over later fields
with its base change, and its rational two-torsion persists.
This proves the exclusion.

For the last assertion, identify the irreducible
four-dimensional Frobenius module with
\(\mathbf F_{\ell^4}\), on which \(F\) multiplies by an
element \(\lambda\).
The symplectic reciprocity involution on its four conjugate
eigenvalues is \(\lambda\mapsto q/\lambda\).
It commutes with their cyclic Galois permutation, is nontrivial,
and has order two. It is therefore the \(\ell^2\)-power map,
giving
\[
                         \lambda^{\ell^2+1}=q.
\]
Since \(q\) generates \(\mathbf F_\ell^\times\),
the cyclic group \(\langle\lambda\rangle\) contains that whole
scalar subgroup. Put \(m=\operatorname{ord}(\lambda)\) and
\(d=m/(\ell-1)\). Then \(\lambda^d\) is a primitive scalar,
while every nonzero vector and quotient functional has exact
period \(m\). The first part applies.
\end{proof}

\paragraph{What the strengthened counts still leave unresolved.}
The preceding statements give a concrete obstruction to treating
the character values in L2 as independent random elements.
Whole primes may fail for representation-theoretic reasons,
before any individual cover is tested. Conversely AC.31 shows
that an ordinary new factor can still be a rational witness even
when Jacobian Frobenius is scalar.

There is also a direct arithmetic limitation on the available
counting criteria. If \(k=\mathbf F_{p^e}\), \(p\) is primitive
modulo \(\ell\), and
\(r=\dim_{\mathbf F_\ell}J_C(k)/\ell J_C(k)>g\), then
\[
                         \ell\le e+1.
\]
Indeed the Frobenius-fixed subspace of \(J_C[\ell]\) is isotropic
for the Weil pairing unless \(q\equiv1\bmod\ell\), and an isotropic
subspace has dimension at most \(g\).
Thus \(r>g\) forces \(\ell\mid q-1\); primitivity gives
\(\ell-1\mid e\).
For large \(e\), a rank greater than \(g\) therefore supplies
only polynomially many cosets, whereas the curve and Jacobian
point counts grow exponentially in \(e\).
In genus two, the threshold in AC.32 consequently still demands
large prime divisors in the lower-rank cases.
No existence theorem for those divisors with the required
residue order has been established here.
The exact Cartier test of AC.27 gives a finite decision for
each proposed character; it does not by itself provide the
uniform torsion-point or Frobenius-orbit avoidance theorem
needed to settle A3 for every curve.

